%% sec-2-background.tex -- Background

\section{Background}

\subsection{3D Gaussian Splatting}
3DGS~\cite{kerbl20233dgs} 将场景表示为一组各向异性 3D 高斯椭球,每个高斯由
位置 / 协方差 / 不透明度 / 球谐系数参数化。渲染阶段使用可微分光栅化器将高斯
投影到屏幕空间,按 alpha-blending 合成像素颜色。

\subsection{4D Gaussian Splatting}
4DGS 在 3DGS 基础上引入时间维度,主流参数化方式分三类:
\begin{itemize}
  \item \textbf{Canonical + Deformation Field}: 4DGS~\cite{wu20244dgs}、
        Deformable-3DGS~\cite{yang2023deformable3dgs}，canonical 空间静态表示
        + per-frame deformation MLP
  \item \textbf{动静态分离}: 4DGS-1K~\cite{yuan20254dgs1k}、MEGA~\cite{zhang2024mega}，
        STV 评分或 buffer-A/B 残差编码
  \item \textbf{Continuous-Time}: RetimeGS~\cite{wang2026retimegs}，消除
        temporal aliasing + ghost-free frame interpolation
\end{itemize}

\subsection{Mobile GPU 基础}
移动端 GPU (Adreno 830) 受限于: 算力 (~2 TFLOPs FP32)、显存带宽、
Vulkan 1.3 驱动成熟度、tile-based rendering 架构对 alpha-blending 的开销。

\bigskip
\noindent\textbf{本节对应 paper notes 索引}: [paper-notes/INDEX.md \S A / B]